\documentclass[11pt,a4paper]{article}

\usepackage[utf8]{inputenc}
\usepackage[T1]{fontenc}
\usepackage[polish]{babel}
\usepackage{lmodern}
\usepackage[a4paper,margin=2.5cm]{geometry}
\usepackage{booktabs}
\usepackage{xcolor}
\usepackage{listings}
\usepackage{tikz}
\usetikzlibrary{positioning,arrows.meta,shapes.geometric}
\usepackage{enumitem}
\usepackage{hyperref}
\hypersetup{colorlinks=true,linkcolor=black,urlcolor=blue,citecolor=black}

% --- listings setup --------------------------------------------------------
\definecolor{codebg}{rgb}{0.97,0.97,0.97}
\definecolor{codecomment}{rgb}{0.40,0.45,0.45}
\definecolor{codekw}{rgb}{0.10,0.10,0.60}
\lstset{
  basicstyle=\ttfamily\footnotesize,
  backgroundcolor=\color{codebg},
  commentstyle=\color{codecomment}\itshape,
  keywordstyle=\color{codekw}\bfseries,
  breaklines=true,
  frame=single,
  rulecolor=\color{gray!40},
  columns=fullflexible,
  showstringspaces=false,
  numbers=left,
  numberstyle=\tiny\color{gray},
  literate={ą}{{\k{a}}}1 {ę}{{\k{e}}}1 {ó}{{\'o}}1 {ł}{{\l{}}}1
           {ś}{{\'s}}1 {ż}{{\.z}}1 {ź}{{\'z}}1 {ć}{{\'c}}1 {ń}{{\'n}}1
}

\title{\textbf{Zarządzanie maszynami wirtualnymi i sieciami\\
je łączącymi na systemie NixOS\\
z wykorzystaniem języka Nix}\\[0.4em]
\large Raport projektowy --- WSO 2026L}
\author{Filip Kobierski (348591) \and Andrzej Pultyn (325213)}
\date{\today}

\begin{document}
\maketitle

\begin{abstract}
\noindent
Celem projektu było opracowanie zestawu narzędzi automatyzujących zarządzanie
maszynami wirtualnymi (VM) oraz \emph{sieciami je łączącymi}. Jako scenariusz
referencyjny przyjęto trójwarstwową usługę WWW (serwer WWW, pamięć podręczna
Redis, baza danych PostgreSQL) wdrożoną jako tzw. \emph{daisy chain} --- łańcuch
maszyn, w którym każde połączenie między warstwami stanowi odrębny segment
sieci. Całość opisano deklaratywnie w języku Nix na systemie NixOS, z
wykorzystaniem KVM/QEMU oraz libvirt. Rozwiązanie zapewnia tworzenie i usuwanie
sieci oraz VM, filtrację ruchu (nftables, polityka \emph{default-deny}), NAT na
brzegu sieci oraz konfigurację zapory dla każdej maszyny. Poprawność i
właściwości bezpieczeństwa potwierdzono zautomatyzowanym testem integracyjnym
realizującym plan testów (widoczność ICMP, skan portów \texttt{nmap}). W
raporcie przeanalizowano bezpieczeństwo rozwiązania, wskazując zarówno wzorce
bezpieczne, jak i typowe nadużycia.
\end{abstract}

\tableofcontents

\section{Scenariusz i zakres}
Zgodnie z tematem projektu zakres obejmuje \emph{zarządzanie sieciami dla
maszyn wirtualnych}: tworzenie i usuwanie sieci oraz VM, filtrację ruchu,
NAT i konfigurację zapory. Scenariusz z koncepcji to działająca usługa
złożona z trzech elementów:
\begin{itemize}[nosep]
  \item serwisu WWW (publiczny punkt wejścia),
  \item pamięci podręcznej Redis,
  \item bazy danych ze zwykłymi danymi (PostgreSQL).
\end{itemize}
Maszyny tworzą \emph{daisy chain}: \texttt{gateway} $\rightarrow$ \texttt{www}
$\rightarrow$ \texttt{cache} $\rightarrow$ \texttt{db}. Struktura ta może w
przyszłości ewoluować do drzewa (np. przez dodanie osobnej, ściślej izolowanej
bazy z danymi krytycznymi) bez zmiany założeń projektu --- wystarczy rozszerzyć
model topologii.

\section{Użyte narzędzia}
\begin{itemize}[nosep]
  \item \textbf{Linux + KVM} --- wirtualizacja sprzętowa,
  \item \textbf{libvirt / QEMU} --- definiowanie sieci i domen (VM),
  \item \textbf{Nix} --- język opisu konfiguracji i budowania artefaktów,
  \item \textbf{NixOS} --- deklaratywny system operacyjny każdej VM,
  \item \textbf{nftables} --- filtracja ruchu i NAT,
  \item \textbf{nginx, Redis, PostgreSQL} --- usługi poszczególnych warstw.
\end{itemize}
Serwerem docelowym jest maszyna z systemem NixOS (\url{http://beka.ovh}).

\section{Architektura}

\subsection{Topologia ,,daisy chain''}
Kluczową decyzją projektową jest umieszczenie \emph{każdego} połączenia między
warstwami w osobnym segmencie warstwy 2 (rys.~\ref{fig:topo}). Dzięki temu nie
istnieje wspólna domena rozgłoszeniowa --- warstwa może komunikować się z inną
tylko wtedy, gdy jest fizycznie podłączona do tego samego segmentu \emph{oraz}
zapora po drugiej stronie na to zezwala.

\begin{figure}[h]
\centering
\begin{tikzpicture}[
  node distance=10mm and 14mm,
  vm/.style={draw,rounded corners,minimum width=20mm,minimum height=12mm,align=center,fill=blue!5},
  cl/.style={draw,rounded corners,minimum width=18mm,minimum height=12mm,align=center,fill=gray!10},
  net/.style={midway,above,font=\scriptsize\ttfamily},
  every edge/.style={draw,thick,-{Latex[length=2mm]}}
]
  \node[cl] (client) {client\\\scriptsize 192.0.2.10};
  \node[vm,right=of client] (gw) {gateway\\\scriptsize NAT/FW};
  \node[vm,right=of gw] (www) {www\\\scriptsize nginx};
  \node[vm,right=of www] (cache) {cache\\\scriptsize redis};
  \node[vm,right=of cache] (db) {db\\\scriptsize pgsql};

  \draw[<->] (client) -- (gw) node[net]{wan};
  \draw[<->] (gw) -- (www) node[net]{edge};
  \draw[<->] (www) -- (cache) node[net]{app};
  \draw[<->] (cache) -- (db) node[net]{data};

  \node[font=\scriptsize,below=2mm of gw] {DNAT 80/443};
  \node[font=\scriptsize,below=2mm of www] {:80,:443};
  \node[font=\scriptsize,below=2mm of cache] {:6379};
  \node[font=\scriptsize,below=2mm of db] {:5432};
\end{tikzpicture}
\caption{Łańcuch maszyn i segmenty sieci je łączące. Klient reprezentuje
dowolny host w Internecie (tylko w teście).}
\label{fig:topo}
\end{figure}

\subsection{Adresacja i macierz dostępu}
Tabela~\ref{tab:nets} przedstawia segmenty i adresy, a tabela~\ref{tab:matrix}
--- jedyne dozwolone przepływy. Wszystko poza nimi jest odrzucane
(\emph{default-deny}).

\begin{table}[h]
\centering
\begin{tabular}{@{}llll@{}}
\toprule
segment & CIDR & VLAN & uczestnicy \\
\midrule
wan  & 192.0.2.0/24 & 9 & gateway \texttt{.1}, client \texttt{.10} (test) \\
edge & 10.10.0.0/24 & 1 & gateway \texttt{.1}, www \texttt{.10} \\
app  & 10.20.0.0/24 & 2 & www \texttt{.1}, cache \texttt{.10} \\
data & 10.30.0.0/24 & 3 & cache \texttt{.1}, db \texttt{.10} \\
\bottomrule
\end{tabular}
\caption{Segmenty sieci i adresacja.}
\label{tab:nets}
\end{table}

\begin{table}[h]
\centering
\begin{tabular}{@{}llll@{}}
\toprule
źródło & cel & port & uzasadnienie \\
\midrule
Internet      & gateway (publ.) & 80, 443 & opublikowany serwis (DNAT $\to$ www) \\
www           & cache           & 6379    & dostęp do Redis \\
cache         & db              & 5432    & dostęp do PostgreSQL \\
zakres admin. & każdy host      & 22      & zarządzanie (SSH) \\
\textit{pozostałe} & --- & --- & \textbf{odrzucane} \\
\bottomrule
\end{tabular}
\caption{Macierz dozwolonego ruchu.}
\label{tab:matrix}
\end{table}

\section{Implementacja}

\subsection{Jedno źródło prawdy: model topologii}
Całe rozwiązanie wyprowadzane jest z jednej struktury danych
(\texttt{lib/topology.nix}). Konfiguracja sieci hostów, reguły nftables, test
integracyjny oraz skrypt zarządzający korzystają z tych samych danych, więc nie
mogą się rozjechać. Listing~\ref{lst:topo} pokazuje fragment opisu hostów.

\begin{lstlisting}[language=,caption={Fragment modelu topologii (lib/topology.nix).},label=lst:topo]
hosts = {
  gateway = { role = "gateway";
    nets = [ { net = "wan";  address = "192.0.2.1"; }
             { net = "edge"; address = "10.10.0.1"; } ]; };
  www = { role = "www";
    nets = [ { net = "edge"; address = "10.10.0.10"; }
             { net = "app";  address = "10.20.0.1";  } ];
    defaultGateway = "10.10.0.1"; };
  cache = { role = "cache";
    nets = [ { net = "app";  address = "10.20.0.10"; }
             { net = "data"; address = "10.30.0.1";  } ]; };
  db = { role = "db";
    nets = [ { net = "data"; address = "10.30.0.10"; } ]; };
};
\end{lstlisting}

Lista \texttt{nets} jest \emph{uporządkowana}: kolejność wpisów wyznacza
kolejność interfejsów sieciowych maszyny. Pozwala to deterministycznie
przypisać sieć logiczną do nazwy interfejsu zarówno w teście NixOS, jak i na
realnym hoście libvirt. Funkcja \texttt{lib/mkHost.nix} zamienia jeden wpis
hosta w kompletny moduł NixOS (adresacja, trasy, zapora, usługi roli).

\subsection{Filtracja ruchu i NAT (nftables)}
Zamiast nadbudowywać reguły nad wbudowaną zaporą NixOS, generujemy
\emph{kompletny}, samodzielny zestaw reguł nftables dla każdej roli
(\texttt{lib/firewall.nix}). W projekcie nastawionym na bezpieczeństwo wartością
jest to, że każdy zaakceptowany pakiet daje się wytłumaczyć na podstawie jednego,
audytowalnego dokumentu. Listing~\ref{lst:db} pokazuje regułę warstwy bazy
danych: ruch przychodzący dozwolony wyłącznie z warstwy cache, a ruch
\emph{wychodzący} domyślnie zabroniony (brak możliwości eksfiltracji).

\begin{lstlisting}[language=,caption={Reguły nftables dla warstwy db (wygenerowane).},label=lst:db]
table inet filter {
  chain input {
    type filter hook input priority 0; policy drop;
    iif "lo" accept
    ct state established,related accept
    ct state invalid drop
    ip protocol icmp icmp type echo-request limit rate 5/second accept
    ip saddr 10.0.0.0/8 tcp dport 22 accept
    ip saddr 10.30.0.1 tcp dport 5432 accept   # tylko z cache
  }
  chain forward { type filter hook forward priority 0; policy drop; }
  chain output {
    type filter hook output priority 0; policy drop;   # brak egress
    oif "lo" accept
    ct state established,related accept
  }
}
\end{lstlisting}

Na bramie (\texttt{gateway}) dodatkowo działa tablica \texttt{nat}: DNAT
publikuje serwis WWW na adresie publicznym, a masquerade umożliwia wybranym
hostom wewnętrznym (warstwie WWW) wyjście do Internetu. Tylko brama ma włączone
przekazywanie pakietów (\texttt{ip\_forward}); pozostałe maszyny są punktami
końcowymi, nie routerami. Warstwy cache i db nie mają trasy domyślnej, więc nie
są w stanie połączyć się z Internetem.

\subsection{Zarządzanie cyklem życia (vmctl)}
Skrypt \texttt{scripts/vmctl.sh} stanowi imperatywne uzupełnienie deklaratywnej
konfiguracji: NixOS decyduje, co znajduje się \emph{wewnątrz} maszyn, a
\texttt{vmctl} tworzy i usuwa \emph{sieci} oraz \emph{domeny} libvirt, na których
te maszyny działają. Segmenty wewnętrzne to sieci izolowane (bez DHCP i bez
routingu hosta --- adresacją zarządza NixOS), a segment \texttt{wan} to sieć NAT
zapewniająca bramie dostęp do Internetu. Najważniejsze podkomendy:
\texttt{up}, \texttt{down}, \texttt{net-up}, \texttt{vm-up}, \texttt{status},
\texttt{console}.

\section{Plan testów i wyniki}
Test integracyjny (\texttt{tests/integration.nix}) uruchamia cały łańcuch w
QEMU --- z tych samych modułów, co wdrożenie produkcyjne --- i realizuje plan
testów z koncepcji, dodając przypadki negatywne, które nadają mu sens:

\begin{itemize}[nosep]
  \item \textbf{Widoczność ICMP:} \texttt{www} pinguje \texttt{cache},
  \texttt{cache} pinguje \texttt{db} (warstwy sąsiednie); \texttt{client}
  \emph{nie} pinguje żadnego hosta wewnętrznego; \texttt{db} \emph{nie}
  inicjuje ruchu (brak egress).
  \item \textbf{Skan \texttt{nmap}:} z Internetu widoczne są wyłącznie porty
  80/443 na adresie publicznym bramy; porty 22, 5432, 6379 są niewidoczne.
  \item \textbf{Osiągalność:} serwis odpowiada przez adres publiczny (DNAT);
  dozwolone są \texttt{www$\to$cache:6379} i \texttt{cache$\to$db:5432};
  zablokowane \texttt{www$\to$db:5432}, \texttt{client$\to$cache:6379},
  \texttt{client$\to$db:5432}.
\end{itemize}

Test uruchamia się poleceniem \texttt{nix flake check} (lub
\texttt{nix build .\#checks.x86\_64-linux.integration -L}) i nie wymaga ręcznej
konfiguracji libvirt.

\section{Analiza bezpieczeństwa}

\subsection{Zastosowane wzorce bezpieczne}
\begin{description}[style=unboxed,leftmargin=1.2em]
  \item[Segmentacja i least privilege.] Jeden segment L2 na połączenie,
  polityka \emph{default-deny}, usługi udostępniane wyłącznie temu sąsiadowi,
  który ich potrzebuje --- z dopasowaniem po \emph{adresie źródłowym}, nie tylko
  porcie. Atakujący, który przejmie warstwę WWW, nie ma drogi sieciowej do bazy
  z pominięciem warstwy cache.
  \item[Obrona w głąb (defense in depth).] Filtracja sieciowa \emph{oraz}
  zabezpieczenia na poziomie usług: Redis nasłuchuje tylko na interfejsie od
  strony \texttt{app} i wymaga hasła; PostgreSQL nasłuchuje tylko na interfejsie
  od strony \texttt{data}, a \texttt{pg\_hba} ogranicza klientów i wymusza
  \texttt{scram-sha-256}.
  \item[Minimalizacja powierzchni ataku.] Brama wystawiona do Internetu nie
  uruchamia żadnej usługi aplikacyjnej. Baza danych nie ma egress ani trasy
  domyślnej, co minimalizuje skutki ewentualnego przejęcia (brak ,,call home'',
  utrudniona eksfiltracja). Wyłączono dokumentację i zbędne pakiety.
  \item[Reprodukowalność jako bezpieczeństwo.] Cała konfiguracja jest
  deklaratywna i wersjonowana, co eliminuje dryf konfiguracji i pozwala
  audytować dopuszczony ruch z jednego modelu. Atomowe aktualizacje i rollback
  NixOS ograniczają ryzyko nieudanych zmian.
  \item[Twardnienie dostępu.] SSH wyłącznie po kluczu (bez haseł, bez logowania
  root), a zapora dodatkowo ogranicza zakres adresów źródłowych mogących w ogóle
  dosięgnąć portu 22.
\end{description}

\subsection{Niewłaściwe sposoby użycia (czego unikać)}
\begin{description}[style=unboxed,leftmargin=1.2em]
  \item[Wystawienie Redis/PostgreSQL na adresie ,,wildcard''.] Redis domyślnie
  nie ma uwierzytelniania, a jego \texttt{protected-mode} łatwo obejść po
  powiązaniu z adresem routowalnym --- otwarty, bezhasłowy Redis to jeden z
  najczęstszych realnych incydentów. Dlatego stosujemy powiązanie z jednym
  interfejsem \emph{i} hasło \emph{i} zaporę.
  \item[Traktowanie NAT jak zapory.] NAT ukrywa topologię, ale nie filtruje
  ruchu w obrębie sieci wewnętrznej. Bezpieczeństwo zapewnia tu nftables, nie
  sam NAT.
  \item[Zbyt szeroki zakres SSH.] Rozszerzenie \texttt{adminCidr} do
  \texttt{0.0.0.0/0} wystawia panel zarządzania na cały Internet.
  \item[Sekrety w repozytorium.] Demonstracyjne \texttt{requirePass} oraz
  poświadczenia bazy są wyłącznie zaślepkami; w produkcji należy dostarczyć je
  poza repozytorium (np. \texttt{requirePassFile}, agenix/sops).
  \item[Spłaszczenie segmentacji.] Połączenie wszystkich warstw w jedną sieć
  ,,dla wygody'' niweczy najważniejszą właściwość bezpieczeństwa --- baza
  staje się osiągalna bezpośrednio z warstwy WWW.
  \item[Nadmierny egress.] Pozostawienie warstwie WWW nieograniczonego wyjścia
  do Internetu (tu: dopuszczone na potrzeby aktualizacji) zwiększa ryzyko
  eksfiltracji; w środowisku produkcyjnym warto ograniczyć je do konkretnych
  celów lub serwera lustrzanego.
\end{description}

\section{Ryzyko i ograniczenia}
Zgodnie z koncepcją głównym ryzykiem była praca z niszowym oprogramowaniem
(ograniczona dokumentacja, mała społeczność). W praktyce najwięcej uwagi
wymagało spójne mapowanie nazw interfejsów: sterownik testów NixOS numeruje
interfejsy pozycyjnie od \texttt{eth1}, podczas gdy realne maszyny libvirt
(z parametrem \texttt{net.ifnames=0}) numerują od \texttt{eth0}. Rozwiązano to
parametrem \texttt{baseIndex} w \texttt{mkHost} oraz dopasowywaniem reguł zapory
po adresach, a nie nazwach interfejsów. Certyfikat TLS warstwy WWW jest
samopodpisany (demonstracyjnie); w produkcji należy użyć ACME lub certyfikatu
dostarczonego.

\section{Wnioski}
Zrealizowano kompletny, deklaratywny zestaw narzędzi do zarządzania maszynami
wirtualnymi i sieciami je łączącymi: tworzenie i usuwanie sieci oraz VM,
filtrację ruchu, NAT i konfigurację zapory dla każdej maszyny. Podejście
oparte na języku Nix dało jedno źródło prawdy, z którego wyprowadzane są
zarówno wdrożenie, jak i automatyczny test potwierdzający właściwości
bezpieczeństwa. Architektura \emph{daisy chain} z odrębnymi segmentami,
polityką \emph{default-deny} i obroną w głąb realizuje zasadę najmniejszych
uprawnień i jest gotowa do rozwoju w strukturę drzewiastą.

\end{document}
